\subsection{Results}
\label{sec:experiments}

For all evaluations, we used proofs generated by Goedel-Prover-V2 \citep{lin2025goedel} on two popular datasets in formal math, miniF2F \citep{zheng2021minif2f} and PutnamBench \citep{tsoukalas2024putnambench}. For miniF2F, we used $n=194$ proofs (average length $334$), and for PutnamBench, we used $n=75$ proofs (average length $1468$).

\subsubsection{Training-Time Results}
In Table~\ref{tab:expert-iteration-results}, we show the min@k and red@k scores of our two models with respect to linted proofs.
We observed steady improvements during each round of expert iteration for both @1 and @32 metrics. \textbf{Our final model outperformed Gemini-2.5-Pro}, a strong reasoning model, even when given proof state annotations similar to Chain-of-States in ImProver~\citep{ahuja2024improver}.


\begin{table}[H]
\centering
\caption{\textbf{Min@k and Red@k throughout expert iteration and online RL.} Our RL model has strong @1 results, while our ExpIt model has strong @32 results. RL metrics are Gaussian-smoothed.}
\label{tab:expert-iteration-results}
\small
\setlength{\tabcolsep}{3pt}
\renewcommand{\arraystretch}{0.95}
\newcommand{\datasetrowlabel}[1]{\raisebox{\dimexpr-3\baselineskip-5pt\relax}[0pt][0pt]{#1}}
\newcommand{\categoryrowlabel}[1]{\raisebox{\dimexpr-\baselineskip\relax}[0pt][0pt]{#1}}
\resizebox{\textwidth}{!}{%
\begin{tabular}{c c c cc cc}
\toprule
\textbf{Dataset} & \textbf{Category} & \textbf{Model} &
\textbf{Min@1 $\downarrow$} & \textbf{Min@32 $\downarrow$} &
\textbf{Red@1 $\uparrow$} & \textbf{Red@32 $\uparrow$} \\
\midrule
\datasetrowlabel{miniF2F}
% ---- ORIGINAL ----
& & \textit{Linted} & \multicolumn{2}{c}{\textit{302}} & \multicolumn{2}{c}{\textit{0.0\%}} \\
\cdashline{2-7}[0.5pt/2pt]

% & \multirow{3}{*}{\centering\textit{Base}}
%   & Qwen2.5-7B  & 294 & 267 & 18.2\% & 35.9\% \\
% &   & Qwen2.5-32B  & 288 & 252 & 23.0\% & 41.8\% \\
&   & Gemini-2.5-Pro & 280 & 207 & 24.3\% & 57.2\% \\
&   & Gemini-2.5-Pro + States & 283 & 207 & 26.4\% & 58.7\% \\
&   & Base (7B)   & 283 & 202 & 17.6\% & 56.2\% \\
\cdashline{2-7}[0.5pt/2pt]
& \categoryrowlabel{\textit{ExpIt}}
  & Base + It 1   & 266 & 178 & 33.4\% & 67.0\% \\
&   & Base + It 2   & 251 & 166 & 45.1\% & 70.6\% \\
&   & \textit{ProofOptimizer-ExpIt}   & 241 & \textbf{153} & 49.0\% & \textbf{72.3\%} \\
\cdashline{2-7}[0.5pt/2pt]
& \textit{RL} & \textit{ProofOptimizer-RL} & \textbf{190} & \textbf{152} & \textbf{63.6\%} & \textbf{70.9\%} \\
  % &  & \textit{It 2 + Test-Time RL} & \textbf{160} & \textbf{154} & \textbf{72.5\%} & \textbf{73.4}\% \\
\midrule
\datasetrowlabel{\shortstack{Putnam\\Bench}}
% ---- ORIGINAL ----
& & \textit{Linted} & \multicolumn{2}{c}{\textit{1359}} & \multicolumn{2}{c}{\textit{0.0\%}} \\
\cdashline{2-7}[0.5pt/2pt]
&  & Gemini-2.5-Pro  & 1348 & 1303 & 5.5\% & 18.0\% \\
&  & Gemini-2.5-Pro + States  & 1371 & 1319 & 6.1\% & 19.2\% \\
&   & Base (7B)   & 1341 & 1222 & 3.9\% & 20.5\% \\
\cdashline{2-7}[0.5pt/2pt]
& \categoryrowlabel{\textit{ExpIt}}
  & Base + It 1   & 1341 & 1215 & 5.2\% & 22.5\% \\
&   & Base + It 2   & 1335 & 1186 & 6.9\% & 24.7\% \\
&   &  \textit{ProofOptimizer-ExpIt}   & 1328 & \textbf{1161} & 8.2\% & \textbf{26.3\%} \\
\cdashline{2-7}[0.5pt/2pt]
& \textit{RL} &  \textit{ProofOptimizer-RL} & \textbf{1303} & 1258 & \textbf{14.9\%} & 21.1\% \\
  % &  & It 2 + Test-Time RL & \textbf{1260} & 1255 & \textbf{23.8\%} & 24.2\% \\
\bottomrule
\end{tabular}
}
\end{table}

% \begin{table}[H]
% \centering
% \begin{tabular}{c c cc cc}
% \toprule
% \textbf{Category} & \textbf{Model} &
% \textbf{Min@1 $\downarrow$} & \textbf{Min@32 $\downarrow$} &
% \textbf{Red@1 $\uparrow$} & \textbf{Red@32 $\uparrow$} \\
% \midrule
% % ---- ORIGINAL ----
% \textit{Linter} & - & \multicolumn{2}{c}{\textit{1359}} & \multicolumn{2}{c}{\textit{0.0\%}} \\
% \noalign{\vskip 1pt}
% \cdashline{1-6}[0.5pt/2pt]
% \noalign{\vskip 2pt}
% % ---- BASE ----
% \multirow{3}{*}{\centering\textit{Base}} & Gemini-2.5-Pro  & 1348 & 1303 & 5.5\% & 18.0\% \\
% & Gemini-2.5-Pro + States  & 1371 & 1319 & 6.1\% & 19.2\% \\
% & Base (7B)   & 1341 & 1222 & 3.9\% & 20.5\% \\
% \noalign{\vskip 1pt}
% \cdashline{1-6}[0.5pt/2pt]
% \noalign{\vskip 2pt}
% % ---- EXPIT ----
% \multirow{3}{*}{\centering\textit{ExpIt}}
%   & Base + It 1   & 1341 & 1215 & 5.2\% & 22.5\% \\
% & Base + It 2   & 1335 & 1186 & 6.9\% & 24.7\% \\
% &  \textit{ProofOptimizer-ExpIt}   & 1328 & \textbf{1161} & 8.2\% & \textbf{26.3\%} \\
% \noalign{\vskip 1pt}
% \cdashline{1-6}[0.5pt/2pt]
% \noalign{\vskip 2pt}
% % ---- GRPO ----
% \multirow{2}{*}{\centering\textit{RL}} &
% \textit{ProofOptimizer-RL} & \textbf{1303} & 1258 & \textbf{14.9\%} & 21.1\% \\
% & It 2 + Test-Time RL & \textbf{1260} & 1255 & \textbf{23.8\%} & 24.2\% \\

% \bottomrule
% \end{tabular}
% \end{table}

% \begin{figure}[H]
%   \centering
%   \begin{minipage}[t]{0.24\textwidth}
%     \centering
%     \includegraphics[width=\linewidth]{figures/grpo/grpo_minif2f_smooth5_small.pdf}
%     \subcaption{miniF2F (train)}
%     \label{fig:rl-training-minif2f}
%   \end{minipage}%
%   \begin{minipage}[t]{0.24\textwidth}
%     \centering
%     \includegraphics[width=\linewidth]{figures/grpo/grpo_putnam_smooth5_nolab_small.pdf}
%         \subcaption{Putnam (train)}
%     \label{fig:rl-training-putnam}
%   \end{minipage}
%   \begin{minipage}[t]{0.24\textwidth}
%     \centering
%     \includegraphics[width=\linewidth]{figures/grpo/ttt_minif2f_nolab_small.pdf}
%     \subcaption{miniF2F (test-time)}
%     \label{fig:ttt-minif2f}
%   \end{minipage}%
%   \begin{minipage}[t]{0.24\textwidth}
%     \centering
%     \includegraphics[width=\linewidth]{figures/grpo/ttt_putnam_nolab_small.pdf}
%         \subcaption{Putnam (test-time)}
%     \label{fig:ttt-putnam}
%   \end{minipage}
%   \caption{\textbf{Evolution of proof reduction (red@1) during RL training (a, b) and test-time RL (c, d).} We use Gaussian smoothing ($\sigma=5$ evaluation intervals for RL training and $\sigma=3$ for test-time RL).}
%   \label{fig:grpo}
% \end{figure}
%
%
% \begin{figure}[htbp]
%   \centering
%   \begin{minipage}[t]{0.4\textwidth}
%     \centering
%     \includegraphics[width=0.9\linewidth]{figures/grpo/ttt_minif2f_small.pdf}
%     \subcaption{miniF2F}
%     \label{fig:ttt-minif2f}
%   \end{minipage}%
%   \hspace{0.1\textwidth}
%   \begin{minipage}[t]{0.4\textwidth}
%     \centering
%     \includegraphics[width=0.9\linewidth]{figures/grpo/ttt_putnam_small.pdf}
%         \subcaption{PutnamBench}
%     \label{fig:ttt-putnam}
%   \end{minipage}
%     \label{fig:ttt}
%     \caption{\textbf{Evolution of relative proof shortening over the course of test-time reinforcement learning.} Shown are periodic evaluation scores for the relative proof shortening metric \textbf{red@1} compared to the original, unlinted proof with Gaussian smoothing with $\sigma=3$ evaluation intervals. The test-time training run jointly uses both datasets as test-time training data and for evaluations.}
% \end{figure}
%
%
% \begin{figure}[htbp]
%     \centering
%     \begin{subfigure}{0.48\textwidth}
%         \centering
%         \includegraphics[width=\linewidth]{figures/grpo/grpo_minif2f_smooth5.pdf}
%         \subcaption{miniF2F}
%         \label{fig:grpo-minif2f}
%     \end{subfigure}
%     \hfill
%     \begin{subfigure}{0.48\textwidth}
%         \centering
%         \includegraphics[width=\linewidth]{figures/grpo/grpo_putnam_smooth5.pdf}
%         \subcaption{PutnamBench}
%         \label{fig:grpo-putnam}
%     \end{subfigure}
%     \caption{\textbf{Evolution of relative proof shortening over the course of reinforcement learning.} Shown are periodic evaluation scores for the relative proof shortening metric \textbf{red@1} with Gaussian smoothing with $\sigma=5$ evaluation intervals. Note that there appears to be a certain tension between metrics on MiniF2F and PutnamBench proofs toward the end of training.}
%     \label{fig:grpo}
% \end{figure}

% \subsubsection{Analysis of Repair with Execution Feedback} \label{subsec:goedel-repair}
% As described in Sec. \ref{subsec:inference-time-algorithms}, we (1) sample $64$ simplifications for each proof with ProofOptimizer-ExpIt, (2) repair incorrect proofs with Goedel-Prover-V2-32B, and (3) shorten successful repairs with our linter. \textbf{Overall, we find while repair with execution feedback leads to improvements, it underperforms resampling because repaired proofs are often even longer than the original proofs.} Fig. \ref{fig:repair-plots} (left) shows the average proof length and reduction \% after sampling, repair, and linting. We find our linter to be effective on repaired proofs, decreasing the average repaired proof length from $644 \to 576$ (miniF2F) and $877 \to 788$ (PutnamBench). In Fig. \ref{fig:repair-plots} (right), we plot the proof length of the original proofs (before Step 1) against simplified proofs (Step 1) and repaired proofs (Step 2). A majority of the repaired proofs (green dots) are above the $y=x$ line, meaning they are longer than the original proofs, let alone the simplified proofs (blue dots).

% \begin{figure}[H]
%     \centering
%     % First plot
%     \begin{subfigure}{0.66\textwidth}
%         \centering
%         \includegraphics[width=\linewidth]{figures/goedel_repair/results_putnam_mini.pdf}
%         % \subcaption{Average Proof Length and Reduction \% for Repair vs. Another Iteration}
%         % \label{fig:repair-results}
%     \end{subfigure}
%     \hfill
%     % Second plot
%     \begin{subfigure}{0.33\textwidth}
%         \centering
%         \includegraphics[width=\linewidth]{figures/goedel_repair/scatter_putnam_mini.pdf}
%         % \subcaption{Comparison of Proof Lengths}
%         % \label{fig:goedel-repair-putnam}
%     \end{subfigure}
%     \caption{Analysis of execution-based repair with Goedel-Prover-V2 on PutnamBench.}
%     \label{fig:repair-plots}
% \end{figure}

% In Table \ref{tab:repair-pipeline-results}, we analyze the success rate of each step of our pipeline. However, the key issue remains to be the high length of the repaired proofs. Even after linting, only 4.8\% (miniF2F) / 1.8\% (Putnam) of post-linted proofs are shorter than the best proof found by ProofOptimizer during simplification.

% \begin{table}
% \caption{Step-by-step success rates, revealing the main bottleneck of long repaired proofs.}
% \label{tab:repair-pipeline-results}
% \centering
% \resizebox{\linewidth}{!}{%
% \begin{tabular}{c c c c }
% \toprule
% \textbf{Dataset} &
% \textbf{\shortstack{Simplification}} &
% \textbf{\shortstack{Repair}} &
% \textbf{\shortstack{Shorter than best (before/after linter)}} \\
% \midrule
% miniF2F &
% $\frac{7852}{12416}$ (63.2\%) &
% $\frac{2840}{4564}$ (62.2\%) &
% $\frac{76}{2840} \to \frac{137}{2840}$ (2.7\% → 4.8\%) \\
% \midrule
% PutnamBench &
% $\frac{1288}{4800}$ (26.8\%) &
% $\frac{613}{3512}$ (17.4\%) &
% $\frac{5}{613} \to \frac{11}{613}$ (0.8\% → 1.8\%) \\
% \bottomrule
% \end{tabular}%
% }
% \end{table}



\subsubsection{Inference-Time Results} \label{subsec:iterative-shortening}

In Fig.~\ref{fig:iterative-shortening} (left), we show the results of iterative proof shortening on miniF2F and PutnamBench proofs using \textit{ProofOptimizer-ExpIt}. First, we did $64$ samples per iteration for $6$ iterations, observing steady improvement at each iteration. To demonstrate the potential of further scaling, we did $1024$ samples at iterations $7$ and $8$ and saw significant improvement. \textbf{Overall, ProofOptimizer combined with iterative proof shortening was very effective on miniF2F and PutnamBench, as average proof length was reduced from $334 \to 75$ and $1468 \to 811$, for an average per-proof reduction of $87.9\%/57.2\%$.} In Fig. ~\ref{fig:iterative-shortening} (right), we plot the overall shortening against the length of the original proof, observing that longer proofs remained challenging to simplify.

\begin{figure}[H]
    \centering
    \begin{subfigure}{0.48\textwidth}
        \centering
        \includegraphics[width=\linewidth]{figures/iterative-simp/results_2x2.pdf}
    \end{subfigure}
    \hfill
    \begin{subfigure}{0.48\textwidth}
        \centering
        \includegraphics[width=\linewidth]{figures/iterative-simp/reduction_vs_length.pdf}
    \end{subfigure}
    \caption{Iterative Shortening: per-iteration improvement (left) and effect of proof length (right)}
    \label{fig:iterative-shortening}
\end{figure}

Finally, in Table~\ref{tab:seed-prover-iterative}, we demonstrate the effectiveness of ProofOptimizer on an out-of-distribution dataset, Seed-Prover's \href{https://github.com/ByteDance-Seed/Seed-Prover/tree/17f89e327e4f90f46b0af385efc233dbbe71f8bb/SeedProver/imo2025/IMO2025}{four IMO 2025 proofs}. With an order of magnitude higher sampling budget, we achieved a significant reduction in the proof length for all four problems, showcasing the potential of our model and technique.

\begin{table}[H]
\caption{Iterative shortening achieves significant reduction for Seed-Prover's IMO 2025 proofs.}
\label{tab:seed-prover-iterative}
\centering
\begin{tabular}{c c c c c}
\toprule
 & \textbf{P1} & \textbf{P3} & \textbf{P4} & \textbf{P5} \\
\midrule
Original Proof Length & 36478 & 16377 & 29147 & 8658 \\
Simplified Proof Length & 20506 & 7907  & 14531 & 4002 \\
Length Reduction & 43.8\% & 51.7\% & 50.1\% & 53.8\% \\
\bottomrule
\end{tabular}
\end{table}

% \begin{table}[H]
% \centering
% \begin{tabular}{c | c | c c c c c c c c c c}
% \toprule
% \textbf{Dataset} & \textbf{Model} & \textbf{Original} & \textbf{Linter} & \textbf{TTT} & \textbf{It 1} & \textbf{It 8} \\
% \midrule
% \multirow{2}{*}{miniF2F}
% & Min@1 & 334 & 302 & 160 & 241 & -- \\
% & Red@1 & 0.0\% & 9.2\% & 75.1\% & 53.9\% & -- \\
% & Min@64 & -- & -- & 152 & 144 & 75 \\
% & Red@64 (\%) & -- & -- & 76.1\% & 76.6\% & 87.9\% \\
% \midrule
% \multirow{2}{*}{Putnam}
% & Min@1 & 1468 & 1359 & 1260 & 1328 & -- \\
% & Red@1 & 0.0\% & 7.4\% & 29.6\% & 15.2\% & -- \\
% & Min@64 & -- & -- & 1253 & 1123 & 811 \\
% & Red@64 (\%) & -- & -- & 30.1\% & 34.8\% & 57.2\%
% \\
% \bottomrule
% \end{tabular}
% \caption{Min@64 (rounded to nearest integer) and reduction (\%) of miniF2F and PutnamBench proofs across inference-time iterations. \red{figure out what to do with this table}}
% % compare to unlinted proof
% \label{tab:inference-iteration-results-horizontal}
% \end{table}
